% GENERATED FILE. Edit canonical lessons, then run npm run generate:books.
% canonical-source-hash: fnv1a64:de5fd07821f95d5c
% canonical-lessons: FA-C454-tamirkar, FA-C454-kaffash, FA-C454-kelidsaz, FA-C454-arayeshgah, FA-C454-saatsaz

\chapter{Repairs around town}
\label{ch:fa-repairs-around-town}
\hlchaptermodality{454}

\begin{chapteropening}
\textbf{By the end of this chapter:} \emph{I can get things fixed in town: the repairman, the cobbler, the locksmith, the hairdresser and the watchmaker.}

Complete the last lesson of chapter 454: I can get things fixed in town: the repairman, the cobbler, the locksmith, the hairdresser and the watchmaker.
\end{chapteropening}

\section[\texorpdfstring{\fa{تعمیرکار} (ta'mirkâr)}{تعمیرکار (ta'mirkâr)}]{\fa{تعمیرکار} (ta'mirkâr) — a repairman, a repairer}

\label{lesson:FA-C454-tamirkar}



\begin{quote}
Before the new one: say the Persian for the traffic police, then the Persian for to ring a number.
\end{quote}

\subsection*{You'll want to know: \fa{تعمیرکار}}
\textbf{\fa{تعمیرکار}} (\emph{ta'mirkâr}) — \textquotedblleft{}a repairman, a repairer\textquotedblright{}.

\begin{quote}
\fa{تعمیرکار} \fa{یخچال} \fa{ما} \fa{را} \fa{درست} \fa{کرد}.
\end{quote}

\begin{quote}
\emph{Ta'mirkâr yakhchâl-e mâ râ dorost kard.}
\end{quote}

\begin{quote}
The repairman fixed our fridge.
\end{quote}

\emph{ta'mir} 'repair' + \emph{kâr} 'worker'.

\begin{cousinweb}
\emph{ta'mirkâr} is built from \emph{ta'mir} 'repair', a loan usually traced to Arabic \emph{taʿmīr}, and native \emph{kâr} 'work, worker'.
\end{cousinweb}

\begin{grammarlens}[title={What you've built}]
One more word for telling, explaining and describing at B1.
\end{grammarlens}

\subsection*{Your turn}
\begin{itemize}
\raggedright
  \item \emph{Say it:} \emph{ta'mirkâr}
  \item \emph{Say it:} \emph{ta'mirkâr}, once more
  \item \emph{Say it:} a sentence of your own with \emph{ta'mirkâr}
  \item \emph{Recall:} say the Persian for the traffic police, then the Persian for to ring a number, then say \emph{ta'mirkâr} again
\end{itemize}

\begin{tcolorbox}[breakable,colback=teal!4,colframe=teal!35!black,title={Before you move on}]
What does \fa{تعمیرکار} mean? (\textquotedblleft{}A repairman, a repairer\textquotedblright{}.) Say it once more.
\end{tcolorbox}
\section[\texorpdfstring{\fa{کفاش} (kaffâsh)}{کفاش (kaffâsh)}]{\fa{کفاش} (kaffâsh) — a shoemaker, a cobbler}

\label{lesson:FA-C454-kaffash}



\begin{quote}
Before the new one: say the Persian for to ring a number, then the Persian for a repairman.
\end{quote}

\subsection*{You'll want to know: \fa{کفاش}}
\textbf{\fa{کفاش}} (\emph{kaffâsh}) — \textquotedblleft{}a shoemaker, a cobbler\textquotedblright{}.

\begin{quote}
\fa{کفاش} \fa{کفش} \fa{پاره} \fa{مرا} \fa{دوخت}.
\end{quote}

\begin{quote}
\emph{Kaffâsh kafsh-e pâre-ye marâ dukht.}
\end{quote}

\begin{quote}
The cobbler stitched my torn shoe.
\end{quote}

Built on \emph{kafsh} 'shoe' in the Arabic pattern of trade names, like \emph{najjâr} 'carpenter'.

\begin{cousinweb}
\emph{kaffâsh} is built on native \emph{kafsh} 'shoe', from Middle Persian \emph{kafš} (EtymDB), shaped in the Arabic pattern of trade names such as \emph{najjâr} 'carpenter'.
\end{cousinweb}

\begin{grammarlens}[title={What you've built}]
One more word for telling, explaining and describing at B1.
\end{grammarlens}

\subsection*{Your turn}
\begin{itemize}
\raggedright
  \item \emph{Say it:} \emph{kaffâsh}
  \item \emph{Say it:} \emph{kaffâsh}, once more
  \item \emph{Say it:} a sentence of your own with \emph{kaffâsh}
  \item \emph{Recall:} say the Persian for to ring a number, then the Persian for a repairman, then say \emph{kaffâsh} again
\end{itemize}

\begin{tcolorbox}[breakable,colback=teal!4,colframe=teal!35!black,title={Before you move on}]
What does \fa{کفاش} mean? (\textquotedblleft{}A shoemaker, a cobbler\textquotedblright{}.) Say it once more.
\end{tcolorbox}
\section[\texorpdfstring{\fa{کلید} \fa{ساز} (kelidsâz)}{کلید ساز (kelidsâz)}]{\fa{کلید} \fa{ساز} (kelidsâz) — a locksmith, a key cutter}

\label{lesson:FA-C454-kelidsaz}



\begin{quote}
Before the new one: say the Persian for a repairman, then the Persian for a shoemaker.
\end{quote}

\subsection*{You'll want to know: \fa{کلید} \fa{ساز}}
\textbf{\fa{کلید} \fa{ساز}} (\emph{kelidsâz}) — \textquotedblleft{}a locksmith, a key cutter\textquotedblright{}.

\begin{quote}
\fa{کلید} \fa{ساز} \fa{برای} \fa{ما} \fa{یک} \fa{کلید} \fa{تازه} \fa{ساخت}.
\end{quote}

\begin{quote}
\emph{Kelidsâz barâ-ye mâ yek kelid-e tâze sâkht.}
\end{quote}

\begin{quote}
The locksmith made a new key for us.
\end{quote}

\emph{kelid} 'key' + \emph{sâz} 'maker'.

\begin{cousinweb}
\emph{kelidsâz} is built from \emph{kelid} 'key', which EtymDB traces through Middle Persian \emph{kilēd} to Greek \emph{kleís} 'key', and native \emph{sâz}, the present stem of \emph{sâkhtan} 'to make'.
\end{cousinweb}

\begin{grammarlens}[title={What you've built}]
One more word for telling, explaining and describing at B1.
\end{grammarlens}

\subsection*{Your turn}
\begin{itemize}
\raggedright
  \item \emph{Say it:} \emph{kelidsâz}
  \item \emph{Say it:} \emph{kelidsâz}, once more
  \item \emph{Say it:} a sentence of your own with \emph{kelidsâz}
  \item \emph{Recall:} say the Persian for a repairman, then the Persian for a shoemaker, then say \emph{kelidsâz} again
\end{itemize}

\begin{tcolorbox}[breakable,colback=teal!4,colframe=teal!35!black,title={Before you move on}]
What does \fa{کلید} \fa{ساز} mean? (\textquotedblleft{}A locksmith, a key cutter\textquotedblright{}.) Say it once more.
\end{tcolorbox}
\section[\texorpdfstring{\fa{آرایشگاه} (ârâyeshgâh)}{آرایشگاه (ârâyeshgâh)}]{\fa{آرایشگاه} (ârâyeshgâh) — a hairdresser's, a barber's (shop)}

\label{lesson:FA-C454-arayeshgah}



\begin{quote}
Before the new one: say the Persian for a shoemaker, then the Persian for a locksmith.
\end{quote}

\subsection*{You'll want to know: \fa{آرایشگاه}}
\textbf{\fa{آرایشگاه}} (\emph{ârâyeshgâh}) — \textquotedblleft{}a hairdresser's, a barber's (shop)\textquotedblright{}.

\begin{quote}
\fa{خواهرم} \fa{پنجشنبه} \fa{به} \fa{آرایشگاه} \fa{رفت}.
\end{quote}

\begin{quote}
\emph{Khâharam panjshanbe be ârâyeshgâh raft.}
\end{quote}

\begin{quote}
My sister went to the hairdresser's on Thursday.
\end{quote}

\emph{ârâyesh} 'adornment' + \emph{gâh} 'place'. The person is an \emph{ârâyeshgar}.

\begin{cousinweb}
\emph{ârâyeshgâh} is a native compound of \emph{ârâyesh} 'adornment', from \emph{ârâstan} 'to adorn', and \emph{gâh} 'place'.
\end{cousinweb}

\begin{grammarlens}[title={What you've built}]
One more word for telling, explaining and describing at B1.
\end{grammarlens}

\subsection*{Your turn}
\begin{itemize}
\raggedright
  \item \emph{Say it:} \emph{ârâyeshgâh}
  \item \emph{Say it:} \emph{ârâyeshgâh}, once more
  \item \emph{Say it:} a sentence of your own with \emph{ârâyeshgâh}
  \item \emph{Recall:} say the Persian for a shoemaker, then the Persian for a locksmith, then say \emph{ârâyeshgâh} again
\end{itemize}

\begin{tcolorbox}[breakable,colback=teal!4,colframe=teal!35!black,title={Before you move on}]
What does \fa{آرایشگاه} mean? (\textquotedblleft{}A hairdresser's, a barber's (shop)\textquotedblright{}.) Say it once more.
\end{tcolorbox}
\section[\texorpdfstring{\fa{ساعت} \fa{ساز} (sâ'atsâz)}{ساعت ساز (sâ'atsâz)}]{\fa{ساعت} \fa{ساز} (sâ'atsâz) — a watchmaker, a watch repairer}

\label{lesson:FA-C454-saatsaz}



\begin{quote}
Before the new one: say the Persian for a locksmith, then the Persian for a hairdresser's.
\end{quote}

\subsection*{You'll want to know: \fa{ساعت} \fa{ساز}}
\textbf{\fa{ساعت} \fa{ساز}} (\emph{sâ'atsâz}) — \textquotedblleft{}a watchmaker, a watch repairer\textquotedblright{}.

\begin{quote}
\fa{ساعت} \fa{پدربزرگ} \fa{را} \fa{پیش} \fa{ساعت} \fa{ساز} \fa{بردیم}.
\end{quote}

\begin{quote}
\emph{Sâ'at-e pedarbozorg râ pish-e sâ'atsâz bordim.}
\end{quote}

\begin{quote}
We took grandfather's watch to the watchmaker.
\end{quote}

\emph{sâ'at} 'watch, hour' + \emph{sâz} 'maker'.

\begin{cousinweb}
\emph{sâ'atsâz} is built from \emph{sâ'at} 'hour, clock', borrowed from Arabic \emph{sāʿa} (EtymDB), and native \emph{sâz}, the present stem of \emph{sâkhtan} 'to make'.
\end{cousinweb}

\begin{grammarlens}[title={What you've built}]
One more word for telling, explaining and describing at B1.
\end{grammarlens}

\subsection*{Your turn}
\begin{itemize}
\raggedright
  \item \emph{Say it:} \emph{sâ'atsâz}
  \item \emph{Say it:} \emph{sâ'atsâz}, once more
  \item \emph{Say it:} a sentence of your own with \emph{sâ'atsâz}
  \item \emph{Recall:} say the Persian for a locksmith, then the Persian for a hairdresser's, then say \emph{sâ'atsâz} again
\end{itemize}

\begin{tcolorbox}[breakable,colback=teal!4,colframe=teal!35!black,title={Before you move on}]
What does \fa{ساعت} \fa{ساز} mean? (\textquotedblleft{}A watchmaker, a watch repairer\textquotedblright{}.) Say it once more.
\end{tcolorbox}
